\documentclass[conference]{IEEEtran}
\usepackage{amsmath,amssymb}
\title{A Queueing View of Tail Latency in Collaborative Inference}
\author{Author A}
\begin{document}
\maketitle

\begin{abstract}
Tail latency in collaborative inference arises from the interaction of queueing, model placement, and cache transfer.
We derive a decomposition that separates the contribution of each stage and validate it on a four-node testbed.
The decomposition explains why adding replicas can increase the 99th percentile when cache transfer saturates the uplink.
\end{abstract}

\section{Introduction}
\label{sec:intro}
Operators often add replicas to reduce tail latency.
In collaborative inference, a new replica also creates new opportunities for cache transfer.
When the uplink is shared, additional replicas can therefore make the tail worse.
We provide a simple decomposition that predicts when this reversal occurs.
The decomposition requires only per-stage service distributions and the transfer volume per request.

\section{Related Work}
\label{sec:related}
Tail-latency studies in datacenters attribute variance to queueing and to background traffic.
Edge studies additionally attribute variance to model heterogeneity and to intermittent connectivity.
Our decomposition combines these factors in a single expression and isolates the transfer term.

\section{Problem Formulation}
\label{sec:problem}
Let $T = T_q + T_s + T_m$ be the request latency, where $T_q$ is queueing delay, $T_s$ is service time, and $T_m$ is migration delay.
Assuming independence across stages, the tail of $T$ is bounded by the sum of stage tails.
The assumption fails when migration and service share the uplink, which we model with a coupling term $\rho$.

\section{Method}
\label{sec:method}
We estimate each stage distribution from traces and compute the coupled tail bound
\begin{equation}
\label{eq:bound}
\mathbb{P}(T > t) \le \sum_{k} \mathbb{P}(T_k > t/3) + \rho(t).
\end{equation}
The coupling term is estimated from concurrent transfer and request traffic.
When $\rho$ grows faster than the queueing term shrinks, adding a replica increases the tail.

\section{Evaluation}
\label{sec:eval}
On a four-node testbed, the bound tracks the measured 99th percentile within 12\%.
Adding a fifth replica reduces the median but increases the 99th percentile by 8\%, as the decomposition predicts.
In particular, the reversal disappears when migration is rate-limited to 20\% of the uplink.

\section{Conclusion}
\label{sec:conclusion}
Tail latency in collaborative inference cannot be reasoned about one stage at a time.
The proposed decomposition identifies when replication backfires.
Future work will extend it to multi-tenant clusters.

\end{document}
